\begin{frame}{자주 하는 실수 (2/2): $P$ 를 ``한 번의 샘플''로 혼동}
  두 번째 흔한 실수: \kb{$P$ 를 ``한 번의 샘플''로 혼동}.
  \begin{itemize}
    \item $P(s'|s,a)$ 는 확률 \textbf{분포}이지 한 번의 결과가
      아니다.
    \item (0에서 right)을 100번 반복하면 \textbf{70번 1로, 30번
      0으로} 갈 것이지만, 100번 중 한 번의 샘플이 ``1로
      갔다''는 사실은 $P = 0.7$ 을 \textbf{증명하지 않는다}
      --- 추정치일 뿐이다.
    \item \textbf{모델 기반} 알고리즘(Week 4--5)은 이 추정치의
      정확도에 \textbf{민감}하므로, 샘플 수와 추정 오차의 관계를
      Week 5에서 자세히 다룬다.
  \end{itemize}
  \vskip0.6em
  \begin{block}{요약}
    ``관측 $\neq$ 상태''(역할의 차이), ``$P$ $\neq$ 한 번의
    샘플''(분포 vs 단일 결과). 두 실수는 모두
    \textbf{확률적 MDP의 표기}와 \textbf{부분관측}을 혼동하는
    것에서 온다.
  \end{block}
\end{frame}
